\chapter{1985年广东卷（理）}

\section{填空题}

\begin{problem}

函数 \(y = \sqrt{\arctan(1 - x)} \) 的定义域是 \fillinblank{}.

\end{problem}

\begin{problem}

如果 \(f\left(\frac{1}{x}\right)=\frac{x}{1-x^{2}} \)，则 \(f(x)= \) \fillinblank{}.

\end{problem}

\begin{problem}

如果一个圆锥底面的半径增大到原来的三倍，但要保持它的体积不变，则其高必须缩小到原来的\fillinblank{}.

\end{problem}

\begin{problem}

在直角坐标系中，通过点 \((0,0) \) \((1,0) \) 和 \((0,2) \) 的圆的方程是\fillinblank{}.

\end{problem}

\begin{problem}

\(\tan\frac{\pi}{5}+\tan\frac{2\pi}{5}+\tan\frac{3\pi}{5}+\tan\frac{4\pi}{5} \) 的值等于 \fillinblank{}.

\end{problem}

\begin{problem}

\(\displaystyle\lim_{n\to\infty}\frac{1+4+7+\cdots+(3n-2)}{n(n+1)} \) 的值等于 \fillinblank{}.

\end{problem}

\begin{problem}

函数 \(y = \sin^2(2x) \) 的最小正周期是 \fillinblank{}.

\end{problem}

\begin{problem}

已知函数 \(f(x) = ax^5 + \arctan x \)，其中 \(a \) 为常数，且 \(f(2) = 10 \)，则 \(f(-2) \) 的值是\fillinblank{}.

\end{problem}

\begin{problem}

已知圆的方程为 \(x^2 + (y - 2)^2 = 9 \)，用平行于 \(x \) 轴的直径把圆分为上下半圆，则以上半圆弧（包括端点）为图象的函数表达式是\fillinblank{}.

\end{problem}

\section{解答题}

\begin{problem}

请把方程 \(y\sqrt{x} - y - 2x + 2\sqrt{x} = 0 \) 所表示的曲线画在右边的直角坐标系中.

\par\noindent
\makebox[\linewidth][r]{\includegraphics[page=1,trim=314.81pt 356.18pt 154.71pt 404.76pt,clip,width=4.2cm]{source/普通高考/1985/1985广东理.pdf}}
\par

\end{problem}

\section{单选题}

\begin{problem}

数 \(y = 2^{-x} \) 的图象的大致形状是

\par\noindent
\makebox[\linewidth][c]{\includegraphics[page=1,trim=31.48pt 147.96pt 78.25pt 597.24pt,clip,width=0.98\linewidth]{source/普通高考/1985/1985广东理.pdf}}
\par

\end{problem}

\begin{problem}

极坐标方程 \(\rho=\frac{10}{5-3\cos\theta} \). 所表示的曲线是

\choiceentry{A}{椭圆}

\choiceentry{B}{圆}

\choiceentry{C}{双曲线}

\choiceentry{D}{抛物线}

\choiceentry{E}{双曲线的一支}

\end{problem}

\begin{problem}

以 \(A(1,3) \)， \(B(-5,1) \) 为端点的线段的垂直平分线的方程是

\choiceentry{A}{\(3x - y + 8 = 0 \)}

\choiceentry{B}{\(3x + y + 4 = 0 \)}

\choiceentry{C}{\(2x - y - 6 = 0\)}

\choiceentry{D}{\(2x + y + 2 = 0 \)}

\choiceentry{E}{\(3x + y + 8 = 0 \)}

\end{problem}

\begin{problem}

如果双曲线 \(\frac{x^{2}}{a^{2}} - \frac{y^{2}}{b^{2}} = 1 \) 的两条渐近线互相垂直，则双曲线的离心率是

\choiceentry{A}{\(\sqrt{2} \)}

\choiceentry{B}{2}

\choiceentry{C}{\(\sqrt{3} \)}

\choiceentry{D}{\(2\sqrt{2} \)}

\choiceentry{E}{\(\frac{3}{2} \)}

\end{problem}

\begin{problem}

设集合 \(C = \{ \) 复数 \(\} \)， \(R = \{ \) 实数 \(\} \)， \(M = \{ \) 纯虚数 \(\} \)，其中 \(C \) 为全集，则

\choiceentry{A}{\(M \cup R = C \)}

\choiceentry{B}{\(\overline{M} \cup R = C \)}

\choiceentry{C}{\(R \cup \overline{R} = C \)}

\choiceentry{D}{\(M \cap R = \{0\} \)}

\choiceentry{E}{\(C \cap \overline{R} = M \)}

\end{problem}

\begin{problem}

\(y = x^{2} (x \leqslant 0) \) 的反函数是

\choiceentry{A}{\(y=\sqrt{x} \)}

\choiceentry{B}{\(y=\pm\sqrt{x} \)}

\choiceentry{C}{\(y=-\sqrt{x} \)}

\choiceentry{D}{\(y = \sqrt{-x} \)}

\choiceentry{E}{\(y = -\sqrt{-x} \)}

\end{problem}

\begin{problem}

如果 \(y = \log_{(a^2 - 1)} x \) 在 \((0, +\infty) \) 内是减函数，则 \(a \) 的取值范围是

\choiceentry{A}{\(|a| > 1 \)}

\choiceentry{B}{\(|a| < \sqrt{2} \)}

\choiceentry{C}{\(a > \sqrt{2} \)}

\choiceentry{D}{\(a < -\sqrt{2} \)}

\choiceentry{E}{\(1 < |a| < \sqrt{2} \)}

\end{problem}

\begin{problem}

如果 \(n\) 是正偶数，则 \(C_{n}^{0} + C_{n}^{2} + \cdots + C_{n}^{n-2} + C_{n}^{n} = \)

\choiceentry{A}{\(2^{n} \)}

\choiceentry{B}{\(2^{n-1} \)}

\choiceentry{C}{\(2^{n-2} \)}

\choiceentry{D}{\((n-1)2^{n-1} \)}

\choiceentry{E}{\((n-1)2^{n-2} \)}

\end{problem}

\begin{problem}

已知 \(180^{\circ} < \alpha < 270^{\circ} \)，且 \(\sin(270^{\circ} + \alpha) = \frac{4}{5} \)，则 \(\tan\frac{\alpha}{2} = \)

\choiceentry{A}{3}

\choiceentry{B}{2}

\choiceentry{C}{\(-2\)}

\choiceentry{D}{\(-3\)}

\choiceentry{E}{\(-\frac{1}{3} \)}

\end{problem}

\begin{problem}

方程 \(\sin 3x - \sin x = 0 \) 的解集是

\choiceentry{A}{\(\{x \mid x = k\pi, k \text{ 是整数}\} \)}

\choiceentry{B}{\(\{x \mid x = \frac{k\pi}{2} + \frac{\pi}{4}, k \text{ 是整数}\} \)}

\choiceentry{C}{\(\{x \mid x = k\pi \text{ 或 } x = \frac{k\pi}{2} + \frac{\pi}{4}, k \text{ 是整数}\} \)}

\choiceentry{D}{\(\{x \mid x = k\pi + \frac{\pi}{2}, k \text{ 是整数}\} \)}

\choiceentry{E}{\(\{x \mid x = \frac{k\pi}{2}, k \text{ 是整数}\} \)}

\end{problem}

\begin{problem}

\(\arcsin\left(-\frac{1}{3}\right),\arctan\left(-\sqrt{2}\right),\arccos\left(-\frac{2}{3}\right)\) 的大小关系是

\choiceentry{A}{\(\arcsin\left(-\frac{1}{3}\right)<\arctan\left(-\sqrt{2}\right)<\arccos\left(-\frac{2}{3}\right) \)}

\choiceentry{B}{\(\arcsin\left(-\frac{1}{3}\right)<\arccos\left(-\frac{2}{3}\right)<\arctan\left(-\sqrt{2}\right) \)}

\choiceentry{C}{\(\arctan\left(-\sqrt{2}\right)<\arcsin\left(-\frac{1}{3}\right)<\arccos\left(-\frac{2}{3}\right) \)}

\choiceentry{D}{\(\arctan\left(-\sqrt{2}\right)<\arccos\left(-\frac{2}{3}\right)<\arcsin\left(-\frac{1}{3}\right) \)}

\choiceentry{E}{\(\arccos\left(-\frac{2}{3}\right)<\arcsin\left(-\frac{1}{3}\right)<\arctan\left(-\sqrt{2}\right) \)}

\end{problem}

\begin{problem}

设向量 \(\overrightarrow{OZ} \) 对应于复数 \(-2\sqrt{3} + 4\i \)，把 \(\overrightarrow{OZ} \) 按顺时针方向旋转 \(60^\circ \) 得到 \(\overrightarrow{OZ_1} \)，则与向量 \(\overrightarrow{OZ_1} \) 对应的复数是

\choiceentry{A}{\(-3\sqrt{3}-\i \)}

\choiceentry{B}{\(\sqrt{3}+5\i \)}

\choiceentry{C}{\(-2\sqrt{3}-4\i \)}

\choiceentry{D}{\(2\sqrt{3} + 4\i \)}

\choiceentry{E}{\(1 + 3\sqrt{3}\i \)}

\end{problem}

\begin{problem}

某班上午要上语文、数学、体育和外语四门课，又体育老师因故不能上第一节和第四节，则不同排课方案的种数是

\choiceentry{A}{24}

\choiceentry{B}{22}

\choiceentry{C}{20}

\choiceentry{D}{12}

\choiceentry{E}{10}

\end{problem}

\begin{problem}

\(\left(x+\frac{1}{\sqrt{x}}\right)^{6} \) 的展开式的常数项是

\choiceentry{A}{1}

\choiceentry{B}{6}

\choiceentry{C}{15}

\choiceentry{D}{20}

\choiceentry{E}{30}

\end{problem}

\begin{problem}

已知 \(\{a_{n}\} \) 为等差数列，且有 \(a_{2} + a_{3} + a_{10} + a_{11} = 48 \)，则 \(a_{6} + a_{7} = \)

\choiceentry{A}{12}

\choiceentry{B}{16}

\choiceentry{C}{20}

\choiceentry{D}{24}

\choiceentry{E}{28}

\end{problem}

\begin{problem}

直角三角形 \(ABC\) 的斜边 \(BC\) 在平面 \(\alpha \) 内，顶点 \(A\) 在平面 \(\alpha \) 外，则 \(\triangle ABC \) 的两条直角边在平面 \(\alpha \) 上的射影与斜边 \(BC\) 所组成的图形只能是

\choiceentry{A}{一条线段}

\choiceentry{B}{一个锐角三角形}

\choiceentry{C}{一个钝角三角形}

\choiceentry{D}{一条线段或一个钝角三角形}

\choiceentry{E}{一条线段或一个锐角三角形}

\end{problem}

\begin{problem}

如果实数 \(x\) 与 \(y\) 满足方程 \(x + y - 4 = 0 \)，则 \(x^{2} + y^{2} \) 的最小值是

\choiceentry{A}{4}

\choiceentry{B}{6}

\choiceentry{C}{8}

\choiceentry{D}{10}

\choiceentry{E}{12}

\end{problem}

\begin{problem}

和 \(y\) 轴相切，且和半圆 \(x^{2} + y^{2} = 4 (0 \leqslant x \leqslant 2) \) 相内切的动圆圆心的轨迹方程是

\choiceentry{A}{\(y^{2} = -4(x - 1)(0 < x \leqslant 1) \)}

\choiceentry{B}{\(y^{2} = 4(x - 1)(0 < x \leqslant 1) \)}

\choiceentry{C}{\(y^{2}=4(x+1)(0<x\leqslant1) \)}

\choiceentry{D}{\(y^{2}=-2(x-1)(0<x\leqslant1) \)}

\choiceentry{E}{\(y^{2}=2(x+1)(0<x\leq1) \)}

\end{problem}

\begin{problem}

设 \(a\)、\(b\)、\(c\) 是空间中的三条直线，下面给出四个命题：

\begin{circlelist}
\item 如果 \(a \perp b\)，\(b \perp c\)，则 \(a \parallel c\).
\item 如果 \(a\)、\(b\) 是异面直线，\(b\)、\(c\) 是异面直线，则 \(a\)、\(c\) 也是异面直线.
\item 如果 \(a\) 和 \(b\) 相交，\(b\) 和 \(c\) 相交，则 \(a\) 和 \(c\) 也相交.
\item 如果 \(a\) 和 \(b\) 共面，\(b\) 和 \(c\) 共面，则 \(a\) 和 \(c\) 也共面.
\end{circlelist}

那么，在上述命题中，真命题的个数是

\choiceentry{A}{4}

\choiceentry{B}{3}

\choiceentry{C}{2}

\choiceentry{D}{1}

\choiceentry{E}{0}

\end{problem}

\begin{problem}

棱柱成为直棱柱的一个必要但不充分的条件是

\choiceentry{A}{棱柱有一条侧棱与底面垂直}

\choiceentry{B}{棱柱有一条侧棱与底面的两条边垂直}

\choiceentry{C}{棱柱有一条侧面与底面的一条边垂直}

\choiceentry{D}{棱柱有一条侧面是矩形，且它与底面垂直}

\choiceentry{E}{棱柱有两个相邻的侧面互相垂直}

\end{problem}

\begin{problem}

用数学归纳法证明： \(|\sin n\theta| \leqslant n|\sin\theta| \).（\(n\) 是自然数）

\end{problem}

\begin{solution}
本题考查绝对值不等式和三角函数的基本知识与数学归纳法的运用.

证： （1）当 \(n=1\) 时，左边 \(|\sin\theta|=\) 右边，所以原不等式成立.

（2）假设当 \(n = k\)（\(k \geq 1\)）时，不等式成立，即
\[
\left|\sin k\theta\right|\leqslant k\left|\sin\theta\right|.
\]
那么，当 \(n=k+1\) 时，
\[
\begin{aligned}
|\sin(k+1)\theta|
&=\left|\sin k\theta\cos\theta+\cos k\theta\sin\theta\right|\\
&\leqslant\left|\sin k\theta\cos\theta\right|+\left|\cos k\theta\sin\theta\right|\\
&=\left|\sin k\theta\right|\left|\cos\theta\right|+\left|\cos k\theta\right|\left|\sin\theta\right|\\
&\leqslant|\sin k\theta|+|\sin\theta|\\
&\leqslant k|\sin\theta|+|\sin\theta|\\
&=(k+1)\left|\sin\theta\right|.
\end{aligned}
\]

这就是说，原不等式对于 \(n = k + 1\) 成立. 根据（1）和（2），原不等式对于任意自然数 \(n\) 都成立.
\end{solution}

\begin{problem}

已知直三棱柱 \(ABC-A_1B_1C_1 \) 的侧棱 \(AA_1=4\text{cm} \)，它的底面 \(\triangle ABC \) 中有 \(AC=BC=2\text{cm} \)，\(\angle C=90^\circ \)，\(E \) 是 \(AB \) 的中点.

\begin{enumerate}
\item 求证：\(CE \perp AB_{1}\).
\item 求证：\(CE\) 和 \(AB_{1}\) 所在的异面直线的距离等于 \(\frac{2\sqrt{3}}{3}\text{ cm}\).
\item 求截面 \(ACB_{1}\) 与侧面 \(ABB_{1}A_{1}\) 所成的二面角的大小（只求其中的锐角，可用反三角函数表示）.
\end{enumerate}

\end{problem}

\begin{solution}
本题考查空间的两条直线的位置关系，二面角大小的计算以及空间想象能力和逻辑推理能力.

（1）证：因为 \(ABC-A_1B_1C_1\) 是直三棱柱，所以 \(AA_{1}\perp\) 底面 \(\triangle ABC\). 又因为 \(CE \subset\) 平面 \(ABC\)，所以 \(CE \perp AA_{1}\). 由于 \(E\) 是等腰三角形 \(ABC\) 底边 \(AB\) 的中点，所以 \(CE \perp AB\). 故 \(CE \perp\) 平面 \(ABB_1A_1\). 因为 \(AB_1 \subset\) 平面 \(ABB_1A_1\)，因此 \(CE \perp AB_1\).

（2）证：因为 \(CE \perp AB\)，\(CE \perp AB_{1}\)，所以 \(CE \perp\) 平面 \(ABB_{1}A_{1}\). 过 \(E\) 作 \(ED \perp AB_1\)，垂足为 \(D\). 因为 \(ED \subset\) 平面 \(ABB_1A_1\)，所以 \(CE \perp ED\). 故 \(ED\) 是 \(CE\) 和 \(AB_{1}\) 所在的异面直线的距离.

在 \(\mathrm{Rt}\triangle ABB_1\) 和 \(\mathrm{Rt}\triangle ADE\) 中，\(\angle A\) 公共，所以 \(\triangle ABB_1 \sim \triangle ADE\). 因此
\[
\frac{ED}{BB_1} = \frac{AE}{AB_1},
\]
即
\[
ED = \frac{AE \cdot BB_1}{AB_1}.
\]
因为 \(AC=BC=2\)，\(\angle ACB=90^{\circ}\)，所以
\[
AB=2\sqrt{2},\quad AE=\frac{1}{2}AB=\sqrt{2}.
\]
又在 \(\mathrm{Rt}\triangle ABB_1\) 中，\(BB_1 = 4\)，所以
\[
AB_1 = \sqrt{AB^2 + BB_1^2} = \sqrt{8 + 16} = 2\sqrt{6}.
\]
因此
\[
ED=\frac{\sqrt{2}\cdot 4}{2\sqrt{6}}=\frac{2\sqrt{3}}{3}\text{ cm}.
\]

（3）因为 \(ED\perp AB_1\)，\(CE\perp\) 侧面 \(ABB_1A_1\)，连结 \(CD\)，依三垂线定理有 \(CD\perp AB_1\). 故 \(\angle CDE\) 是截面 \(CAB_1\) 和侧面 \(ABB_1A_1\) 所成的二面角的平面角，且为锐角. 又 \(CE\) 是 \(\mathrm{Rt}\triangle ABC\) 的斜边上的中线，所以 \(CE=\frac{1}{2}AB=\sqrt{2}\). 由上可知 \(ED=\frac{2\sqrt{3}}{3}\). 因此在 \(\mathrm{Rt}\triangle CED\) 中，
\[
\tan\angle CDE=\frac{CE}{ED}=\frac{\sqrt{6}}{2}.
\]
故截面 \(CAB_{1}\) 和侧面 \(ABB_{1}A_{1}\) 所成的较小的二面角等于 \(\arctan\frac{\sqrt{6}}{2}\).
\end{solution}

\begin{problem}

设函数 \(f(x)=\log_{3}\left(x^{2}-4mx+4m^{2}+m+\frac{1}{m-1}\right) \)，其中 \(m\) 是实数，又用 \(M\) 表示集合 \(\{m \mid m>1\} \).

\begin{enumerate}
\item 求证：当 \(m \in M \) 时，\(f(x)\) 对所有实数 \(x\) 都有意义；反之，如果 \(f(x)\) 对所有实数 \(x\) 都有意义，则 \(m \in M\).
\item 当 \(m \in M \) 时，求函数 \(f(x)\) 的最小值.
\item 求证：对每一个 \(m \in M \)，函数 \(f(x)\) 的最小值都不小于 \(1\).
\end{enumerate}

\end{problem}

\begin{solution}
本题考查二次函数与对数函数的基本性质，不等式的证明和不等式的解法.

（1）证：当 \(m \in M\) 时，有 \(m > 1\)，从而对所有的实数 \(x\) 都有
\examdisplaymath{x^{2}-4mx+4m^{2}+m+\frac{1}{m-1}=(x-2m)^{2}+m+\frac{1}{m-1}\geqslant m+\frac{1}{m-1}>0.}
于是，当 \(m \in M\) 时，函数 \(f(x)\) 对所有的实数 \(x\) 都有意义.

反之，如果 \(f(x)\) 对所有的实数 \(x\) 都有意义，则对所有的实数 \(x\)，都有
\examdisplaymath{x^{2}-4mx+4m^{2}+m+\frac{1}{m-1}>0.}
而当 \(x=2m\) 时，上式变为
\examdisplaymath{m+\frac{1}{m-1}>0,}
即
\examdisplaymath{\frac{m^{2}-m+1}{m-1}>0}
或
\examdisplaymath{\frac{\left(m-\frac{1}{2}\right)^{2}+\frac{3}{4}}{m-1}>0.}
由于上式左端的分子总是正数，所以它的分母 \(m-1>0\)，即 \(m>1\)，从而 \(m\in M\).

（2）因为以 \(3\) 为底的对数函数是增函数，所以
\examdisplaymath{f(x)=\log_{3}\left[(x-2m)^{2}+m+\frac{1}{m-1}\right]\geqslant\log_{3}\left(m+\frac{1}{m-1}\right).}
又因为
\examdisplaymath{f(2m)=\log_{3}\left(m+\frac{1}{m-1}\right),}
所以，当 \(m \in M\) 时，\(f(x)\) 的最小值是
\[
\log_{3}\left(m+\frac{1}{m-1}\right).
\]

（3）证：当 \(m \in M\) 时，有 \(m > 1\)，所以
\[
m+\frac{1}{m-1}=(m-1)+\frac{1}{m-1}+1\geqslant2\sqrt{(m-1)\frac{1}{m-1}}+1=3.
\]
上式两端取以 \(3\) 为底的对数，得
\[
\log_{3}\left(m+\frac{1}{m-1}\right)\geqslant\log_{3}3=1.
\]
于是，根据（2）的结果可知，对每一个 \(m \in M\)，函数 \(f(x)\) 的最小值都不小于 \(1\).
\end{solution}

\begin{problem}

在直角坐标系中，有椭圆 \(I\) 和抛物线 \(II\)，它们的参数方程分别是 \(I\)：\(\left\{\begin{array}{l} x = m + 2\cos\phi,\\ y = \sqrt{3}\sin\phi \end{array}\right.\)（\(m\) 是常数，\(\phi\) 是参数）；\(II\)：\(\left\{\begin{array}{l} x = \frac{3}{2}+t^{2},\\ y = \sqrt{6}t \end{array}\right.\)（\(t\) 是参数）.

\begin{enumerate}
\item 求证：当 \(m=4\) 时，椭圆 \(I\) 有一个焦点和一条准线分别与抛物线 \(II\) 的焦点和准线重合.
\item 求证：当且仅当 \(m \in [-\frac{1}{2}, \frac{7}{2}]\) 时，椭圆 \(I\) 与抛物线 \(II\) 有交点.
\item 在区间 \([-\frac{1}{2}, \frac{7}{2}]\) 中有哪些 \(m\) 值，使得椭圆 \(I\) 和抛物线 \(II\) 的交点与坐标原点的距离等于这个交点与椭圆 \(I\) 的中心的距离？
\end{enumerate}

\end{problem}

\begin{solution}
本题考查二次曲线的基础知识，曲线参数方程的应用，以及逻辑推理能力和计算能力.

（1）证：将曲线 \(I\)，\(II\) 的方程化为直角坐标方程：
\[
I:\ \frac{(x-m)^{2}}{4}+\frac{y^{2}}{3}=1,\qquad II:\ y^{2}=6\left(x-\frac{3}{2}\right).
\]
当 \(m = 4\) 时，椭圆 \(I\) 中心的坐标为 \((4,0)\)，长轴在 \(x\) 轴上，长半轴长 \(a = \sqrt{4} = 2\)，短半轴长 \(b = \sqrt{3}\). 从而半焦距
\[
c = \sqrt{a^2 - b^2} = 1,
\]
焦点坐标 \((4 \pm c,0)\) 即为 \((5,0)\) 和 \((3,0)\)，准线方程
\[
x = 4 \pm \frac{a^2}{c},
\]
即为 \(x = 8\) 和 \(x = 0\).

由抛物线 \(II\) 的方程可知：该抛物线以 \(x\) 轴为对称轴，开口向右，顶点的坐标为 \(\left(\frac{3}{2},0\right)\)，焦点至顶点的距离是 \(\frac{p}{2}=\frac{3}{2}\). 从而焦点的坐标为
\[
\left(\frac{3}{2}+\frac{p}{2},0\right)=(3,0),
\]
准线方程是
\[
x=\frac{3}{2}-\frac{p}{2}=0.
\]
所以，椭圆 \(I\) 的左焦点和左准线分别与抛物线 \(II\) 的焦点和准线重合.

（2）证：将曲线 \(I\)，\(II\) 的直角坐标方程联立，得
\[
\left\{
\begin{aligned}
\frac{(x-m)^{2}}{4}+y^{2}&=1,\quad \circled{1}\\
y^{2}&=6\left(x-\frac{3}{2}\right).\quad \circled{2}
\end{aligned}
\right.
\]
那么，曲线 \(I\)，\(II\) 有交点，当且仅当方程组有实数解 \(x\)、\(y\). 将 \circled{2} 代入 \circled{1}，整理得
\examdisplaymath{x^{2}+(8-2m)x+m^{2}-16=0.\quad \circled{3}}
又由于
\[
x - \frac{3}{2} = \frac{y^2}{6} \geq 0,
\]
故方程组有实数解，当且仅当方程 \circled{3} 有不小于 \(\frac{3}{2}\) 的实根.

方程 \circled{3} 的两个根为
\[
x_{1}=-(4-m)+2\sqrt{2(4-m)},\qquad x_{2}=-(4-m)-2\sqrt{2(4-m)}.
\]
可见，当且仅当 \(4 - m \geqslant 0\) 时，\(x_1\)、\(x_2\) 是实根，且这时 \(x_2 \leqslant 0\). 由此可知：方程 \circled{3} 有不小于 \(\frac{3}{2}\) 的实根，当且仅当 \(m\) 满足不等式组
\[
\begin{cases}
4-m\geqslant0,\\
-4+m+2\sqrt{2(4-m)}\geqslant\frac{3}{2}.
\end{cases}
\]
解此不等式组的解集是 \(\left\{m \mid -\frac{1}{2} \leqslant m \leqslant \frac{7}{2}\right\}\). 综合起来，即得：当且仅当 \(m \in [-\frac{1}{2}, \frac{7}{2}]\) 时，椭圆 \(I\) 与抛物线 \(II\) 有交点.

（3）因为椭圆 \(I\) 中心的坐标为 \((m,0)\)，所以曲线 \(I\)，\(II\) 的交点 \((x, y)\) 满足题设，等价于 \(x\)、\(y\) 满足
\[
x^{2} + y^{2} = (x - m)^{2} + y^{2},
\]
即
\[
m(m-2x)=0.\quad \circled{1}
\]
由（2）的证明可知，曲线 \(I\)，\(II\) 的交点横坐标是
\[
x = (m - 4) + 2 \sqrt{2(4 - m)},\qquad m \in \left[-\frac{1}{2}, \frac{7}{2}\right].
\]
将它代入 \circled{1} 式，得知所求的 \(m\) 值应满足方程
\[
m \left[ 8 - m - 4 \sqrt{2 (4 - m)} \right] = 0,\qquad m \in \left[-\frac{1}{2}, \frac{7}{2}\right].
\]
解方程：由 \(m = 0\) 得 \(m_{1} = 0\)；由 \(8 - m - 4\sqrt{2(4 - m)} = 0\)，即
\examdisplaymath{m^{2} + 16m - 64 = 0,}
得 \(m_{2} = 8(\sqrt{2} -1)\in \left[-\frac{1}{2},\frac{7}{2}\right]\)，\(m_{3} = -8(\sqrt{2} +1)\notin \left[-\frac{1}{2},\frac{7}{2}\right]\)（应舍去）. 直接检验可知，当 \(m = 0\) 或 \(8(\sqrt{2} - 1)\) 时，曲线 \(I\)，\(II\) 的交点满足题设条件，故所求的 \(m\) 的值是 \(0\) 和 \(8(\sqrt{2} - 1)\).
\end{solution}

\additionalproblemsection

\begin{problem}

已知函数 \(f(x)=x^{5}-5x+c \)，其中 \(c\) 是实数.

\begin{enumerate}
\item 求 \(f(x)\) 的极大值和极小值.
\item 证明方程 \(f(x)=0\) 的不同实根的个数不大于 \(3\).
\end{enumerate}

\end{problem}

\begin{solution}
本题考查利用导数求函数极值和讨论方程的不同的实根个数的能力.

（1）由 \(f'(x)=5x^{4}-5=5(x^{4}-1)\). 令 \(f'(x)=0\)，即 \(5(x^{4}-1)=0\)，这个方程的实根是 \(x=\pm1\). 而
\[
f(-1) = -1 + 5 + c = c + 4,\qquad f(1) = 1 - 5 + c = c - 4.
\]
函数 \(f(x)\) 在 \((-\infty,-1)\) 与 \((1,+\infty)\) 上递增，在 \((-1,1)\) 上递减. 所以 \(f(x)\) 的极大值是 \(f(-1)=c+4\)，\(f(x)\) 的极小值是 \(f(1)=c-4\).

（2）证：用反证法证明. 假如方程 \(f(x) = 0\) 的不同实根个数大于 \(3\)，则至少可找到四个不同的实数 \(x_{1} < x_{2} < x_{3} < x_{4}\)，使得
\[
f(x_{1}) = f(x_{2}) = f(x_{3}) = f(x_{4}) = 0.
\]
根据微分中值定理，有 \(\theta_{1}\in(x_{1},x_{2})\)，使得
\[
f(x_{2}) - f(x_{1}) = f'(\theta_{1})(x_{2} - x_{1}).
\]
由上式及 \(x_{2} > x_{1}\)，可知 \(f'(\theta_1) = 0\). 这表明方程 \(f'(x) = 0\) 在区间 \((x_{1},x_{2})\) 中至少有一个实根 \(\theta_{1}\). 同理可证：方程 \(f'(x)=0\) 在区间 \((x_{2}, x_{3})\)、\((x_{3}, x_{4})\) 中分别有实根 \(\theta_{2}\)、\(\theta_{3}\). 显然，有 \(\theta_{1}<\theta_{2}<\theta_{3}\)，即方程 \(f'(x)=0\) 至少有三个不同的实根，这与（1）中所得的结果相矛盾. 从而，方程 \(f(x)=0\) 的不同实根的个数不大于 \(3\).
\end{solution}

\begin{problem}

本题考查基础知识和基本运算，只需要直接写出结果.

\end{problem}

\begin{problem}

\(\{x \mid x \leq 1\} \) 2. \(\frac{x}{x^{2}-1} \) 3. \(\frac{1}{9} \)

\end{problem}

\begin{problem}

\(x^{2} + y^{2} - x - 2y = 0 \) 5. 0 6. \(\frac{3}{2} \) 7. \(\frac{\pi}{2} \)

\end{problem}

\begin{problem}

-10 9. \(y=2+\sqrt{9-x^{2}}(-3\leqslant x\leqslant3) \) 10. 如图

\end{problem}

\begin{problem}

本题考查基础知识和基本运算. 1. E 2. A 3. B 4. A 5. C 6. C 7. E 8. B 9. D 10. C 11. C 12. B 13. D 14. C 15. D 16. D 17. C 18. A 19. E 20. B

\end{problem}
